\documentclass[10pt,a4paper]{amsart}
\usepackage[margin=19mm]{geometry}
\usepackage{amsmath,amssymb,mathtools}
\usepackage[scr=rsfs,cal=euler]{mathalfa}
\usepackage{xfrac}
\usepackage[unicode,hidelinks,bookmarks=false]{hyperref}
\newcommand{\ie}{i.e.\ }
\newcommand{\cf}{cf.\ }
\newcommand{\resp}{respectively\ }
\newcommand{\Subsection}{\subsection}
\providecommand{\nobreakdash}{\nobreak\hbox{-}}
\setlength{\emergencystretch}{2em}
\multlinegap0pt
\usepackage{bm}
\usepackage{bbm}
\usepackage{dsfont}
\usepackage{xspace}
\usepackage{cancel}
\usepackage{mathtools}
\usepackage{amsthm}
\makeatletter
\@ifundefined{theorem}{\newtheorem{theorem}{Theorem}}{}
\@ifundefined{lemma}{\newtheorem{lemma}[theorem]{Lemma}}{}
\@ifundefined{proposition}{\newtheorem{proposition}[theorem]{Proposition}}{}
\@ifundefined{remark}{\newtheorem{remark}[theorem]{Remark}}{}
\makeatother
\DeclareMathOperator{\Gal}{Gal}
\DeclareMathOperator{\NS}{NS}
\newcommand{\MW}{\mbox{MW}}
\newcommand{\Z}{\mathbb{Z}}
\newcommand{\ra}{\rightarrow}
\title[Rank jumps and Multisections of elliptic fibrations on K3 su]{Rank jumps and Multisections of elliptic fibrations on K3 surfaces}
\author{Alice Garbagnati, Cec\'ilia Salgado}
\date{Source: arXiv:2505.15159v2 (CC BY 4.0)}
\begin{document}
\maketitle

\noindent\emph{Source and licence.} Rank jumps and Multisections of elliptic fibrations on K3 surfaces. Version arXiv:2505.15159v2 (CC BY 4.0), distributed under the Creative Commons Attribution 4.0 International licence, as recorded in the arXiv licence metadata for this exact version and shown on the abstract page https://arxiv.org/abs/2505.15159v2. Full rights notice: licenses/arxiv-2505.15159-CC-BY-4.0.txt.\par\smallskip

\noindent\textit{Editorial scope.} Counted unit: the Theorem whose source label is \texttt{thm: rank jump odd d} in the article (source0.tex lines 695--710, source label \texttt{thm: rank jump odd d}), delivered together with the article's own complete original proof of it (source0.tex lines 712--728, one \texttt{proof} environment of 17 lines). No mathematics is written, repaired or re-derived: statement and proof are the article's own text line for line. Required context retained uncounted: lem:indep\_multisections (lines 283--285); theorem:Picard3 (lines 355--374); eq: explcit of bisection (lines 551--551); rem; intersection B and sections (lines 552--556); prop: Fratx cover of P2 (lines 568--574). Every definition, display and label of those lines is retained unchanged. Omitted from this package and not redistributed: 1--282, 286--354, 375--550, 557--567, 575--694, 711--711, 729--1050. Nothing inside a retained excerpt is omitted or altered: every retained body is delivered exactly as the article's lines are written, so no omission receipt of any kind accompanies this package. Citations: the retained excerpts carry no citation command at all, so no bibliography entry of the article is transcribed and no reference is redistributed. Reference adaptation: none was needed and none was made; every reference inside the delivered unit resolves inside it, so no unresolved reference or citation is produced. Printed slips kept literal: no printed word, symbol, hypothesis or number of the counted statement or of its proof was altered. Layout and class: the article's own class is \texttt{amsart} with packages including fullpage, inputenc, amssymb, amsmath, enumerate, epsfig, color, caption, longtable, graphics, multirow, ulem, graphicx, xypic; the wrapper is the lane's pinned \texttt{amsart} editorial wrapper at 10pt on A4, which restates the theorem-like environments the retained text needs and adds the article's own macro definitions that the retained text needs (\texttt{\textbackslash Gal}, \texttt{\textbackslash MW}, \texttt{\textbackslash NS}, \texttt{\textbackslash Z}, \texttt{\textbackslash ra}), with extra wrapper packages (bm, bbm, dsfont, xspace, cancel, mathtools). The delivered numbering is the unit's own. Assets: no retained span contains a figure environment, a graphics-inclusion command, a PDF-page inclusion, a table or a TikZ picture; no image file is copied into this package, the package declares no asset dependency and no image file is redistributed with it. Licence: version arXiv:2505.15159v2 of this article is distributed under the Creative Commons Attribution 4.0 International licence, as recorded in the arXiv licence metadata for this exact version and shown on the abstract page https://arxiv.org/abs/2505.15159v2. A full rights notice is delivered at licenses/arxiv-2505.15159-CC-BY-4.0.txt. Characters: every retained excerpt is pure ASCII, so the delivered engine, XeLaTeX with its default Latin Modern text font, reports no missing character.
\begin{lemma}\label{lem:indep_multisections}
Let $\pi:X \rightarrow C$ be an elliptic fibration and $M$ a saliently ramified multisection defined over a field $l$. Let $X_M$ and $\pi_M: X_M \rightarrow M$ be as above. Then the section $\sigma_M$ induced on $X_M$ by $(\iota, \mathrm{id}):M \rightarrow X \times_C M$ is linearly independent of the pull-back of the sections of $\pi$ in $\MW(X_M,\pi_M,l)$.
\end{lemma} 

\begin{theorem}\label{theorem:Picard3}
Let $X$ be a K3 surface with $\rho(X)=3$ and $X\in\mathcal{L}_d$. 

If $d=1$, then the unique elliptic fibration on $X$ has a reducible fiber, its Mordell--Weil group is trivial and there are no smooth rational or genus 1 bisections. 

For any other $d$, any elliptic fibration on $X$ has no reducible fibers, the Mordell--Weil group of any elliptic fibration on $X$ is isomorphic to $\Z$ and admits a generator that intersects the zero section in $d-2$ points. 

Moreover, the following hold:

\begin{itemize}
\item[a)] If $d\leq 3$, the unique elliptic fibration on $X$ does not admit smooth multisections of genus $g\leq 1$.

\item[b)] If $d>3$, then
\begin{itemize}
\item
any elliptic fibration on $X$ admits a smooth rational bisection if, and only if, $d\equiv 1\mod 2$ and this bisection is salient. 
\item any elliptic fibration on $X$ admits a smooth genus 1 bisection if, and only if, $d\equiv 0 \mod 2$ and this bisection is salient. 
\end{itemize}
\end{itemize}
\end{theorem}

\begin{equation}\label{eq: explcit of bisection}B=\frac{d+3}{2}F+2\mathcal{O}+b_3.\end{equation}

\begin{remark}\label{rem; intersection B and sections}
In the previous situation (i.e. $\NS(X)\simeq \Lambda_d$, $d\equiv 1\mod 2$, $d\geq 5$), we have shown that every elliptic fibration on $X$ admits a smooth rational bisection $B$ such that $B\mathcal{O}=x=\frac{d-5}{2}$, whose class is given in \eqref{eq: explcit of bisection}. The parity of $B\mathcal{O}$ is the same as the one of $BS_n$ for every $n\in\mathbb{Z}$ and depends on $d$. More specifically, $BS_n=\frac{d+3}{2}+2dn(n-1)-4\equiv \frac{d+3}{2}\mod 4$, i.e. $BS_n\equiv x\mod 4$. In particular:  
\begin{eqnarray}\label{eq: intersection B sections}\begin{array}{c}B\mathcal{O}\equiv BS_n\equiv 0\mod 2 \Leftrightarrow d\equiv 1\mod 4\\
B\mathcal{O}\equiv BS_n\equiv 1\mod 2 \Leftrightarrow d\equiv 3\mod 4.\end{array}\end{eqnarray}
\end{remark}

\begin{proposition}\label{prop: Fratx cover of P2}
Let $X\in\mathcal{L}_{2x+5}$, then there exists $f:X\ra \mathbb{P}^2$ a double cover of $\mathbb{P}^2$ branched on a plane sextic $C_6$ with a simple node at a point $p$, such that there exists a rational curve $R\subset\mathbb{P}^2$ of degree $x+1$ passing through $p$ with multiplicity $x$ that splits in the double cover.

The pencil of lines through $p$ induces an elliptic fibration on $X$ whose zero section is one of the components of $f^{-1}(R)$.

The cover involution of $f:X\ra\mathbb{P}^2$ is the involution $ T_1\circ \varepsilon$ and preserves the class of the bisection $B$.
\end{proposition}

\begin{theorem}\label{thm: rank jump odd d}
Let $(X, \pi)$ be a generic elliptic K3 surface with Picard number 3, i.e. $X$ is generic in a family $\mathcal{L}_d$.  Let $k$ be the field of definition of $\pi$, $d$ odd and $d>3$. 
\begin{itemize}
\item 
If $d\equiv 1\mod 4$, then there is an extension $l/k$ of degree at most two such that the set 
\[
\mathcal{R}(X,\pi, l)=\{t\in \mathbb{P}^1(l); \mathrm{rank} X_t(l)>\mathrm{rank} \MW(X,\pi,l)\}
\] is infinite. 
In particular, $X$ has the potential rank jump property.
\item If $d\equiv 3\mod 4$ then the set 
\[
\mathcal{R}(X,\pi, k)=\{t\in \mathbb{P}^1(k); \mathrm{rank} X_t(k)> \mathrm{rank} \MW(X,\pi,k)\}
\] is infinite. 
In particular,   $X$ has the rank jump property. 
\end{itemize}
\end{theorem}

\begin{proof}
By Theorem \ref{theorem:Picard3}, under the assumption that $d>3$ is odd,  for a generic $X$  the fibration $\pi$ admits a smooth salient bisection $B$ that is of genus 0. By Lemma \ref{lem:indep_multisections}, $B$ yields a section that is linearly independent from the sections induced by $\pi$ on the base change.

Therefore, the only statement which has not yet been proved concerns the field of definition $k$ (resp. $l$), namely, to show that the sets $\mathcal{R}(X,\pi, k)$ (resp. $\mathcal{R}(X,\pi, l)$) are infinite if $d\equiv 3\mod 4$ (resp. $d\equiv 1\mod 4$). In particular will  verify that $B$ is defined over $k$ if $d$ is odd, that it has infinitely many $k$-points if $d\equiv 3\mod 4$ and that it has infinitely many $l$-points for a certain field extension $l$ of degree two if $d\equiv 1\mod 4$. 

\textbf{Claim 1:} The field of definition of the rational bisection from Theorem \ref{theorem:Picard3} is $k$. 

Indeed, under the assumption that $d$ is odd, a generic K3 surface in a family $\mathcal{L}_d$ with Picard number 3 admits a realization as the double cover of $\mathbb{P}^2_{k}$ branched on a sextic with a unique node (Proposition \ref{prop: Fratx cover of P2}). The genus 1 fibration is defined by the pull-back of the pencil of lines through the node and thus the branch sextic and its node are defined over the same field $k$. The exceptional divisor above the node is isomorphic to a rational curve defined over $k$ and is a bisection of $\pi$. In particular, the bisection is defined over the same field as $\pi$. 

In particular, since $X$ is a double cover of the blow up of $\mathbb{P}^2_{k}$, and $B$ is defined over $k$, it follows that $B$ has infinitely many $l$-points for a certain degree two extension $l/k$.

\textbf{Claim 2:} The bisection $B$ from Theorem \ref{theorem:Picard3} has infinitely many points over $k$ if $d\equiv 3\mod 4$.

Since $B$ is rational and defined over $k$, it suffices to show that there is at least one point of $B$ defined over $k$ to conclude that there are infinitely many.

The Galois group $\Gal(l/k)$ acts on the points of $B$ by permuting the ones that are only defined over $l$ (and not over $k$) and fixing the ones defined on $k$. If $d\equiv 3 \mod 4$, then there is an odd number of points in $B\mathcal{O}$, by Remark \ref{rem; intersection B and sections}. Both the curves $B$ and $\mathcal{O}$ are defined over $k$, hence the Galois action preserves both of them. In particular it preserves the set of their intersection points. Since $\Gal(l/k)$ has order 2, if it acts on an odd number of points, then it fixes at least one of them, so there is a point in $B\cap \mathcal{O}$ (and hence on $B$) which is defined over $k$. Therefore $B$ can be embedded as a conic in the plane. The presence of a $k$-point allows us to take the projection away from it and conclude that $B$ is isomorphic to $\mathbb{P}^1$ over $k$.
\end{proof}

\end{document}
